\begin{table}[ht]
\caption{\textbf{\econ example}}
\centering
% \resizebox{\textwidth}{!}{
\begin{tabular}{p{13cm}}
\toprule
% \hspace{5.5cm} \econ \\
% \midrule
\textit{\textbf{Query}} \\
\midrule
Would a GDP measure be improved by excluding foreign interest paid? \\
The income method of calculating GDP is as follows: GDP = wages + profits + rents + interest + depreciation + taxes + NFFI. If an economy has high external debt, for instance, because it used external financing to buy machinery and equipment, then foreign interest payments will be high. In that case, wouldn't GDP (per capita) be a poor measure of economic well-being since a significant portion of the generated income is leaving the home economy? \\
Wouldn't excluding foreign interest payments in calculating GDP/ income give us a better measure since this represents what income is available to the home economy? \\
\midrule
\textit{\textbf{Reasoning steps}} \\
\midrule
The statement implies that GDP serves as an indicator of economic well-being, yet GDP, Gross domestic product,  primarily reflects economic production and outputs rather than directly gauging well-being. It's crucial to understand the association between GDP and economic well-being. \\
\midrule
\textit{\textbf{Example positive document}} \\
\midrule
Measuring well-being is about going beyond the cold numbers of GDP. But cold as it may be, GDP remains a very important indicator for measuring the economic performance of countries, which is a fundamental driver of well-being. This is not only because income means higher material standards but also because income is needed to sustain private and public spending in non-material components of well-being, such as education and health. However, there is an increasing recognition that focusing on GDP alone is not enough to achieve better lives for all. \\
The relationship between a country’s s GDP per capita and well-being is usually positive – in other words, countries with higher GDP per capita are also those where well-being is higher on average. However this relationship becomes weaker as a country’s income grows, suggesting that once income reaches a certain level, increased income is less likely to generate well-being. The other interesting phenomenon is that some countries do better at delivering wellbeing than they do if gauged only on the basis of economic production per capita. This is the case for Nordic European countries but also for New Zealand. On the other hand,  there are countries that do better in GDP per capita than on average well-being,  for instance the United States and Switzerland. \\
% So why in these two sets of countries do economic performance and well-being  not go strictly  hand in hand? One explanation is that the countries that do better in terms of well-being have  made the choice of working less to achieve a better work and life balance. This translates into lower income but also increased leisure time that can be shared with friends and family, or that is used for volunteering and engaging with the community. Another reason is that these countries have better environmental quality, partly because of lower economic production and thus pollution. Well-being is a matter of democracy and societal choices, what really matters at the end of the day is whether countries are where their citizens would like them to be. It can be higher economic production or more time for life, longer healthy lives or more connected neighbourhoods. What is important is to collect the relevant set of statistics to judge whether collective objectives are met and ensure that governments are paying the right degree of attention to these statistics.\\
... \\
\midrule
\textit{\textbf{Example negative document}} \\
\midrule
The impact of foreign interest rates on the economy: The role of the exchange rate regime\\
It is often argued that many economies are affected by conditions in foreign countries. This paper explores the connection
between interest rates in major industrial countries and annual real output growth in other countries. The results show that high
foreign interest rates have a contractionary effect on annual real GDP growth in the domestic economy, but that this effect is
centered on countries with fixed exchange rates. The paper then examines the potential channels through which major-country
interest rates affect other economies. The effect of foreign interest rates on domestic interest rates is the most likely channel when
compared with other possibilities, such as a trade effect. \\
... \\
This paper answers two questions. First, what is the effect of interest rates in base countries on other countries' annual
real GDP growth?3 Second, how does this effect vary by the exchange rate regime and other country characteristics?
Answering the second question helps to disentangle the channels through which foreign country interest rates affect
other economies. We find that annual real output growth in countries is negatively associated with interest rates in their
base countries, but that this effect holds only for countries with fixed exchange rates. \\
... \\
\bottomrule
\end{tabular}
% }
\label{tab:economic_example}
\end{table}